% ======================================================================
% 03-method.tex — the AI-assisted reusability loop
% ======================================================================
\section{Method: the AI-assisted reusability loop}
\label{sec:method}

The project did not begin with a software-engineering plan. It began with a
research question and a set of runs whose results did not survive scrutiny. The
method emerged as a response to that: a six-step loop, executed by a human and
an agent, with the repository as the memory between turns.

\subsection{Roles}
\label{sec:roles}

\begin{table}[htbp]
\centering
\small
\begin{tabularx}{\textwidth}{@{}lX@{}}
\toprule
Actor & Responsibility in this project \\
\midrule
Human (PI/editor) &
Chooses the question and the acceptance threshold; decides which evidence
class a claim belongs to; rejects the best-seed number; writes and owns the
final prose; accepts or rejects a policy. \\
Pi (agent harness) &
Reads the repository; runs commands and tests; edits files; writes new tests,
harnesses and documentation; commits with the finding in the message; keeps the
memory files current. \\
DeepSeek (model) &
Proposes the diagnosis, the patch, the test and the summary that the harness
executes. Its output is a hypothesis until a gate passes. \\
Gemini (tertiary) &
Occasional research assistance and code review; not part of the loop. \\
\bottomrule
\end{tabularx}
\caption{Division of labour. The harness/model split matters: most of the
leverage came from the agent's ability to \emph{run} the check, not from a
single model completion.}
\label{tab:roles}
\end{table}

\subsection{The loop}
\label{sec:loop}

Each turn has six steps (Figure~\ref{fig:loop}):
\begin{enumerate}
  \item \textbf{Read the memory.} Before touching code, read \code{HUMAN.md},
        \code{Actions.md}, \code{Conclusions.md} and the relevant plan. This is
        the step that prevents re-learning an environment gotcha or
        re-litigating a closed question.
  \item \textbf{Reproduce the claim.} Run the command that produced the number,
        or the smallest version of it. A claim that cannot be reproduced is a
        bug report, not a baseline.
  \item \textbf{Change one thing.} The smallest reversible edit that could
        explain the failure---not a cleanup, not a refactor.
  \item \textbf{Run the gate.} A named threshold on a named artifact. If the
        gate does not exist yet, writing it is part of the change.
  \item \textbf{Record the finding.} Update the run log and the conclusions,
        including negative results. A finding that is not written where the next
        session will look did not happen.
  \item \textbf{Commit the finding.} The commit message states what was learned,
        not what files moved. This makes the history searchable by claim.
\end{enumerate}

\begin{figure}[htbp]
\centering
\begin{tikzpicture}[
  node distance=7mm and 9mm,
  box/.style={draw, rounded corners, align=center, font=\small,
              minimum height=9mm, text width=24mm, fill=blue!4},
  mem/.style={draw, dashed, rounded corners, align=center, font=\scriptsize,
              minimum height=8mm, text width=40mm, fill=orange!8},
  arrow/.style={-{Latex[length=2mm]}, thick}
]
\node[box] (read) {1. Read\\memory};
\node[box, right=of read] (repro) {2. Reproduce\\the claim};
\node[box, right=of repro] (change) {3. Smallest\\change};
\node[box, below=of change] (gate) {4. Run the\\gate};
\node[box, left=of gate] (record) {5. Record\\the finding};
\node[box, left=of record] (commit) {6. Commit\\the finding};
\node[mem, below=of record] (memory) {repository as memory:\\docs, manifests, tests, artifacts};
\draw[arrow] (read) -- (repro);
\draw[arrow] (repro) -- (change);
\draw[arrow] (change) -- (gate);
\draw[arrow] (gate) -- (record);
\draw[arrow] (record) -- (commit);
\draw[arrow] (commit.north) -- ++(0,4mm) -| (read.south);
\draw[arrow, dashed] (memory) -- (record);
\draw[arrow, dashed] (memory) -- (read);
\end{tikzpicture}
\caption{The loop. The dashed arrow is the property that makes it cumulative:
the repository, not the conversation, is the memory.}
\label{fig:loop}
\end{figure}

The loop is deliberately unoriginal; it is the scientific method with a commit
hook. What changed with AI is the \emph{cost per turn}. A human will run a
five-seed sweep once, at the end, when the paper is due. An agent will run it on
every candidate because the marginal cost of ``run the same command four more
times and tabulate the result'' is small. That is the entire mechanism behind
the negative result in \S\ref{sec:failures}: the $\tau$-ramp experiment was repeated
because repeating it was cheap, and the repetition destroyed the effect.

\subsection{Artifacts as memory}
\label{sec:memory}

Conversation context is volatile; a repository is not. The project therefore
keeps four kinds of memory, each written for a different reader.

\begin{table}[htbp]
\centering
\small
\begin{tabularx}{\textwidth}{@{}lX@{}}
\toprule
Artifact & Question it answers \\
\midrule
\code{HUMAN.md} & ``What will bite me when I come back?'' Environment gotchas,
wrong interpreter, out-of-memory, buffering, the corner artefact. Written for a
human under time pressure. \\
\code{Actions.md} & ``What was run, and what happened?'' A chronological run log
with numbers, artifact paths and negative results. \\
\code{Conclusions.md} & ``What is trustworthy, what is single-seed, what is
open?'' The evidence-class ledger for the whole project. \\
\code{USAGE.md} & ``How do I run this?'' Every entry point, its cost and its
outputs. \\
\code{runs/runtimes.jsonl} & ``What did it cost?'' One JSON record per run, with
git SHA, grid, seed, device, wall clock and status. \\
\code{out/*/run\_config.json} & ``What exactly was this?'' The effective config,
seed, device, dtype and metrics next to the model. \\
Commit messages & ``Why does this line exist?'' Each message states the finding,
so \code{git log -{}-grep} is a literature search over the project's own history. \\
\bottomrule
\end{tabularx}
\caption{The memory hierarchy. The first four are prose for humans; the last
three are machine-readable records for reruns and audits.}
\label{tab:memory}
\end{table}

Two rules keep the memory useful. First, \textbf{a finding is written at the
level of the reader who will need it}: a run goes in \code{Actions.md}, a
conclusion goes in \code{Conclusions.md}, an environment trap goes in
\code{HUMAN.md}. Second, \textbf{numbers carry their provenance}: a table entry
names the artifact directory and the harness, so ``median $1.61\times10^{-2}$''
can be traced to \code{runs/sweeps/baseline\_n5} and re-derived. The article's
own \code{EVIDENCE.md} is the same idea applied to the manuscript.

\subsection{What the agent actually did}
\label{sec:agent-work}

Across the repository's history at the time of writing: 69 commits, 304 tracked
files, 23{,}739 lines of Python, 22 test modules and 172 test functions, eight
optimization configurations, and 33 runtime records (Figure~\ref{fig:timeline}).
The agent's work was not uniformly distributed. It was concentrated in exactly
the activities that are repetitive and easy to defer:
\begin{itemize}
  \item writing the parity tests that compare the NumPy reference and the torch
        backend scenario by scenario;
  \item turning a debugging script into a harness (\code{seed\_sweep.py}), then
        using it to overturn a single-seed result;
  \item adding a guard after every failure (a trivial-predictor baseline, a
        non-finite-candidate policy, a zip-on-crash path);
  \item rebuilding documentation tables from artifacts rather than memory;
  \item reconstructing a timeline of 69 commits to find when a claim changed.
\end{itemize}

\begin{figure}[htbp]
\centering
\begin{tikzpicture}[x=1cm,y=1cm,font=\scriptsize]
  \draw[-{Latex[length=2mm]}, thick] (0,0) -- (12.7,0);
  \foreach \x/\d in {0/11-13, 3/11-16, 6/11-18, 9/11-19, 12/11-21} {
    \draw (\x,0.12) -- (\x,-0.12);
    \node[below] at (\x,-0.16) {\d};
  }
  \node[above, align=center, text width=2.8cm] at (0,0.35)
    {initial commit\\(PINN/FNO/opt tree)};
  \node[above, align=center, text width=3.0cm] at (3,0.95)
    {seed variance;\\($\tau$-ramp negative)};
  \node[above, align=center, text width=3.0cm] at (6,0.35)
    {FNO coordinate fix;\\Conclusions.md};
  \node[above, align=center, text width=3.0cm] at (9,0.95)
    {backend parity;\\Colab/fleet drivers};
  \node[above, align=center, text width=3.2cm] at (12,0.35)
    {certificates;\\42-page assembly};
  \draw[gray, dashed] (0,0) -- (0,0.25);
  \draw[gray, dashed] (3,0) -- (3,0.85);
  \draw[gray, dashed] (6,0) -- (6,0.25);
  \draw[gray, dashed] (9,0) -- (9,0.85);
  \draw[gray, dashed] (12,0) -- (12,0.25);
\end{tikzpicture}
\caption{The eight-day span in which the loop ran. The dates are commit dates, not a
claim about effort; the point is the density of findings per day once the gates
and the harness existed.}
\label{fig:timeline}
\end{figure}
Those are not the tasks that make a research career, which is precisely why
they are the ones that get skipped. The rest of the article describes what they
bought.

\subsection{Threats to the method}
\label{sec:method-threats}

The loop is not self-validating. An agent can write a test that encodes the
current bug, a gate with a threshold chosen to pass, or a summary that rounds a
single seed into a trend. The defences used here were procedural, not
architectural: the human set thresholds before seeing results; every merged
mechanism had to fail on a known-bad input; and negative results were recorded
with the same care as positive ones. These defences are fallible, and
\S\ref{sec:discussion} returns to them.
